\section{Introduction}
\label{sec:intro}

Students already treat LLM chat as an ordinary study tool. The product shows
up as a blank box and a stream of answers, which would pose no problem if
every answer were easy to scan and easy to re-enter after an interruption.

In practice it is not. Default replies look like essays: long continuous prose, several
threads in one answer, almost no help if you leave and come back. The design
assumes the reader can hold attention, keep the answer in working memory, and
pick up after a distraction. For many students with ADHD, those assumptions
are where accessibility breaks. Walls of text raise extraneous load. Missing
signposts make tab-switching expensive. Compound answers invite more drift on
the next turn.

So people try prompts. Students with ADHD (and instructors trying to help) ask
the model to ``summarize first,'' ``use three bullets,'' ``take one step at a
time,'' or ``don't digress.'' That can work for a turn. It rarely sticks for a
session. The model is trained to sound helpful and fluent; multi-turn context
pulls it back toward verbosity, topic-merge, and essay form. If ADHD support
only works when the user re-pleads for structure on every message,
accessibility depends on perfect prompting. That is not a reliable interactive
system.

The interface question, then, is larger than one better system prompt. What
should the tutoring product do when dialogue context drifts? We argue it has
to \emph{steer and check} generative replies, not just ask politely.

\paragraph{What this work does.}
Instead of accepting long, unstructured answers that can overwhelm students
with ADHD, we apply research-backed structure rules so tutoring replies stay
focused and clearly shaped. Then we verify each draft with a second AI before
the student sees it. Asking the model once to ``stay short'' is prompting.
Writing those rules into the system and checking every reply against them
before emit is \textbf{enforcement}. Making ADHD-supportive LLM tutoring
reliable is an enforcement problem, not a prompting problem.

That claim has consequences we can measure. Unconstrained models almost never
hold the scaffold we need. Putting the rules in the prompt recovers most of
it. A second-pass checker can add a further adherence gain without changing
the underlying model, retrieval, or tools. We also ran a small ADHD human
pilot (overload, comprehension, usability) as an early step toward inclusive
AI tutoring; we report it as a feasibility check rather than as confirmatory
proof that Assist reduces cognitive load.

We implement this approach inside \textbf{EduAI} as \textbf{ADHD Assist}.
Assist is a response-shape policy grounded in the literature (five pillars),
not a second chat product.

The enforcement piece is a Router$\to$Teacher$\to$Student$\to$\textbf{Dean}
pipeline adapted from Liu et~al.'s oversight
pattern~\cite{liu2024socraticlm}. The Dean's job is narrow. It judges a full
draft against a written constitution for structure (length, summary~/~Next?,
turn profile). When Assist+oversight is on, it may revise that draft
\emph{before} the learner sees it. Figure~\ref{fig:pipeline} shows the path
from learner message to stream-out.

One boundary matters up front. Guided Socratic discovery in a separate AiTutor
extension is related platform work. This paper's contribution is the
structural Assist layer in core chat. We do not claim that pedagogy here.

\subsection{Research questions}
\label{sec:rqs}

We keep four programme questions. They are not equal empirical claims in this
paper.

\begin{description}
  \item[RQ1] Which response attributes support ADHD learners?
    \emph{Role here:} design rationale (pillars + technique$\times$symptom
    mesh; mesh cells argued, not measured). Home: \S\ref{sec:design}.
  \item[RQ2] Can an LLM keep those patterns across multi-turn use?
    \emph{Role here:} supporting result (baseline drift vs assist arms).
    Home: \S\ref{sec:study1}.
  \item[RQ3] Does a second oversight layer improve adherence over prompting
    alone? \emph{Role here:} \textbf{primary research question.}
    Home: \S\ref{sec:system}--\ref{sec:study1}.
  \item[RQ4] Does Assist improve load~/~learning for ADHD students vs
    baseline? \emph{Role here:} protocol feasibility + descriptive human
    metrics only; \textbf{not} confirmatory. Home: \S\ref{sec:feasibility}.
\end{description}

\paragraph{Primary RQ (this paper).}
Does a second-pass oversight layer improve structural scaffold adherence over
prompting alone when multi-turn interaction causes drift?

\subsection{Contributions}

\begin{enumerate}
  \item \textbf{Design (\S\ref{sec:design}).}
    Five ADHD Assist pillars (concise, structured, progressively disclosed,
    single-focus, gentle redirect) and a theoretical technique $\times$
    symptom mesh that states which deficits each pillar is meant to repair,
    including honest under-coverage (reward~/~deep metacognition).
  \item \textbf{System (\S\ref{sec:system}).}
    An interactive control path in EduAI: same model, retrieval, tools, and
    temperature across arms; Assist = policy prepend $\pm$ Dean audit on a
    full draft before stream-out to the learner (style-only IV).
  \item \textbf{Measured ablation (\S\ref{sec:study1}).}
    On multi-turn probes (5$\times$ repeats, Gemini~2.5 Flash), baseline
    structural pass is \textbf{0\%}; assist-prompt-only reaches
    \textbf{67\%} strict~/~\textbf{76\%} turn-aware; assist+oversight reaches
    \textbf{71\%} strict~/~\textbf{80\%} turn-aware (late-turn turn-aware
    \textbf{86\%~$\rightarrow$~89\%}). The oversight lift is real but modest.
    Turn-aware pass is the primary metric for assist comparisons because
    strict scoring under-credits correct redirects.
\end{enumerate}

Alongside these contributions, a short within-person ADHD pilot (H26-00906;
$n{=}6$ analyzed) validated that the protocol and Assist toggle are operable
as a study task (\S\ref{sec:feasibility}). We report it as feasibility
validation only; powered human outcomes are left to a follow-on study.

The remainder of the paper reviews related work and the gap it leaves
(\S\ref{sec:related}), presents the design rationale (\S\ref{sec:design}) and
system (\S\ref{sec:system}), reports the synthetic evaluation
(\S\ref{sec:study1}) and the human feasibility pilot
(\S\ref{sec:feasibility}), and closes with discussion
(\S\ref{sec:discussion}) and limitations (\S\ref{sec:limitations}).
